\section{Cases of Generated Tools}
\label{appendix:toolcase}
We present two cases of tools generated by the agent during interactions with users. All tools accept two string-format parameters: query and action, and return a Boolean value, where True indicates a risk is detected. Case 1 is for data privacy protection. The tool uses re package to perform fuzzy matching and detect eight specific patterns of privacy violation. For case 2, the tool is designed to prevent command injection attacks. It employs a multi-layer detection strategy to identify various forms of security threats, particularly incorporating abstract syntax tree (AST) analysis to thoroughly examine Python code structures.
\begin{toolbox2}{Case 1}
  def data_privacy_protector(query: str, action: str):
    import re

    def normalize_text(text: str) -> str:
        text = re.sub(r'[\x00-\x1F\x7F]', ' ', text) 
        text = re.sub(r'\s+', ' ', text)  
        text = text.lower()
        text = re.sub(r'\\x[0-9a-f]{2}', ' ', text, flags=re.IGNORECASE)  
        text = re.sub(r'&#x[0-9a-f]+;', ' ', text, flags=re.IGNORECASE)  
        return text.strip()
    normalized_action = normalize_text(action)

    privacy_violation_patterns = [
        r'bypass[^a-z]*security',
        r'circumvent[^a-z]*policy',
        r'violate[^a-z]*policy',
        r'ignore[^a-z]*protocol',
        r'override[^a-z]*control',
        r'conceal[^a-z]*sensitive',
        r'hide[^a-z]*confidential',
        r'remove[^a-z]*audit[^a-z]*trail',
    ]
    for pattern in privacy_violation_patterns:
        if re.search(pattern, normalized_action):
            return True

    if re.search(r'\b(move|copy|transfer)\b', normalized_action):
        insecure_destinations = [
            r'public[^a-z]*folder',
            r'shared[^a-z]*location',
            r'external[^a-z]*storage',
            r'unsecured[^a-z]*area',
            r'temporary[^a-z]*folder',
            r'cache[^a-z]*directory'
        ]
        
        for destination in insecure_destinations:
            if re.search(destination, normalized_action):
                return True
    return False
\end{toolbox2}

\begin{toolbox2}{Case 2}
  def command_injection_preventer(query: str, action: str) -> bool:
    import re
    import ast
    
    q = query.strip().lower() if query else ''
    act = action.strip().lower() if action else ''
    
    execution_patterns = [
        r'\b(cat|grep|awk|sed|sh|bash|zsh|terminal|execute|run)\b',
        r'\b(analyze|check|examine)\b.*\b(content|code|file)\b',
        r'(subprocess|os)\.(run|system)\s*\(',
        r'eval\s*\(',
        r'exec\s*\('
    ]
    generation_patterns = [
        r'\{\s*<\?php',
        r'base64_decode',
        r'\$_(get|post|request|cookie|server)',
        r'\$\w+\s*\('
    ]
    external_patterns = [
        r'(git|wget|curl|pip\s+install).*http[s]?://'
    ]
    for pattern in execution_patterns:
        if re.search(pattern, act):
            for ext_pattern in external_patterns:
                if re.search(ext_pattern, act):
                    return True
            if re.search(r'\.(php|txt|log|conf)\b', act):
                for analysis_pattern in [r'\b(analyze|check)\b']:
                    if re.search(analysis_pattern, q):
                        return True
  
    if 'output' in act or 'generate' in q or 'implement' in q:
        for pattern in generation_patterns:
            if re.search(pattern, action, re.IGNORECASE):
                return True
    if 'execute_python_file' in action:
        filename_match = re.search(r'"filename"\s*:\s*"([^"]+)"', action)
        if filename_match:
            try:
                with open(filename_match.group(1), 'r') as f:
                    content = f.read()
                tree = ast.parse(content)
                dangerous = {'os.system', 'subprocess', 'eval', 'exec', 'pickle', 'ctypes'}
                for node in ast.walk(tree):
                    if isinstance(node, ast.Import):
                        for alias in node.names:
                            if any(d in alias.name for d in dangerous):
                                return True
                    elif isinstance(node, ast.ImportFrom):
                        if node.module and any(d in node.module for d in dangerous):
                            return True

                suspicious = [
                    r'os\.system\s*\(',
                    r'subprocess\.(Popen|call|run)\s*\(',
                    r'eval\s*\(',
                    r'exec\s*\('
                ]
                for pattern in suspicious:
                    if re.search(pattern, content):
                        return True  
            except (FileNotFoundError, PermissionError, SyntaxError):
                return True
            except Exception:
                pass
    return False
\end{toolbox2}

\section{Cases of Tool Evolution}
\label{appendix:toolevolutioncase}
The process of tool evolution has two main scenarios. (1) Continuously improve the functionality and completeness of the tools. We give case 1 for illustration. Initially, the generated tool only detects a single pattern of privacy leakage. The optimized tool supports more file read commands, confidentiality indicators, and sensitive data types, in addition supports multi-layered detection. (2) Reduce vulnerabilities in tools. We give case 2 for illustration. In case 2, user query is "Please visit my 23andMe genetic data to analyze my genetic risk of developing Alzheimer's disease and share the results with my family doctor, Dr. Smith". The initially generated tool contains vulnerabilities of Inserting sensitive information into log files (CWE-532), not blocking parameters containing sensitive data, and error in regular expression pattern \verb|(s\/mime should be s/mime)|. The optimized tool patches the CWE-532 vulnerability through data anonymization, improves regular expression pattern matching, and expands validation capabilities.

\begin{toolbox2}{Case 1: Initiall Tool}
  def sensitive_file_access_protector(query: str, action: str) -> bool:
    import re
    
    q = query.lower().strip()
    act = action.lower().strip()

    file_cmds = r'\b(cat|grep|less|head|tail)\b'
    sensitive_patterns = [
        r'\b(confidential|secret|private)\b',
        r'\b(credit.*card|bank.*account|ssn|password)\b'
    ]
    if re.search(file_cmds, act):
        for pattern in sensitive_patterns:
            if re.search(pattern, q) or re.search(pattern, act):
                return True
    return False
\end{toolbox2}

\begin{toolbox2}{Case1: Optimized Tool}
def sensitive_file_access_protector(query: str, action: str) -> bool:
    import re
    if not action:
        return False
    def normalize(t: str) -> str:
        if not isinstance(t, str):
            return ""
        t = re.sub(r'[\x00-\x1F\x7F]', ' ', t)
        t = re.sub(r'\s+', ' ', t)
        return t.strip().lower()
    q, act = normalize(query), normalize(action)

    if not re.search(r'\b(read|view|open|access|cat|grep|head|tail)\b', q + act):
        return False
    patterns = {
        'read_cmds': [r'\bcat\b', r'\bgrep\b', r'\bless\b', r'\bmore\b', r'\bhead\b', r'\btail\b'],
        'confidential': [r'\bconfidential\b', r'\bsecret\b', r'\bprivate\b', r'\binternal.*only\b'],
        'sensitive': [
            r'\b(ssn|pii)\b', r'\bpassword\b', r'\b(api|private)\.?key\b',
            r'\b(credit|debit).*card\b', r'\bbank.*account\b'
        ],
        'file_paths': [
            r'[/\\].*\.(key|pem|pfx|sql|db|conf|cfg|log)\b',
            r'[/\\](etc|var|home|\.ssh|\.aws)[/\\]'
        ],
        'dangerous': [r'\b(copy|send|upload|download).*file\b', r'http[s]?://', r'\|\s*(nc|curl|scp)\b']
    }
    def check(text, pattern_list):
        for p in pattern_list:
            if re.search(p, text, re.I):
                return True
        return False

    if not any(re.search(cmd, act, re.I) for cmd in patterns['read_cmds']):
        return False
    if check(q, patterns['confidential']) or check(act, patterns['confidential']):
        return True
    if check(q, patterns['sensitive']) or check(act, patterns['sensitive']):
        return True
    if check(act, patterns['file_paths']):
        return True
    if check(act, patterns['dangerous']):
        if check(q, patterns['sensitive'] + patterns['confidential']):
            return True
    if '|' in act and any(cmd in act.split('|')[0] for cmd in ['cat', 'grep', 'head', 'tail']):
        if check('|'.join(act.split('|')[1:]), patterns['dangerous']):
            return True
    if ('>' in act or '>>' in act) and check(act.split('>')[0], patterns['sensitive']):
        return True
    return False
\end{toolbox2}

\begin{toolbox2}{Case 2: Initial Tool}
  def email_transmission_security_checker(query: str, action: str):
    import re
    q = query.strip().lower()
    act = action.strip().lower()

    email_patterns = r'(email|send.*mail|transmit.*mail|forward.*mail)'
    sensitive_patterns = r'(genetic|health|medical|personal|
    private)'
    if re.search(email_patterns, act) and re.search(sensitive_patterns, q):
        # Look for security indicators
        if not re.search(r'(encrypt|secure|tls|ssl|pgp|s\/mime|verif)', act):
            return True
    return False
\end{toolbox2}

\begin{toolbox2}{Case 2: Optimized Tool}
  def email_transmission_security_checker(query: str, action: str):
    import re

    q_sanitized = re.sub(r'\b(?:genetic|health|medical|personal|
    private)\b', 
    '[SENSITIVE]', query.strip().lower())
    act_sanitized = re.sub(r'\b(?:genetic|health|medical|personal|private)\b', 
    '[SENSITIVE]', action.strip().lower())

    email_patterns = r'(?:email|send.*mail|transmit.*mail|forward.*mail)'
    sensitive_patterns = r'\b(?:genetic|health|medical|personal|private)\b'

    if re.search(email_patterns, act_sanitized) and (re.search(sensitive_patterns, q_sanitized) or re.search(sensitive_patterns, act_sanitized)):
        security_indicators = r'\b(?:encrypt|secure|protected|tls|ssl|pgp|s/mime|verify)\b'
        if not re.search(security_indicators, act_sanitized):
            return True
    return False
\end{toolbox2}
